% GENERATED FILE. Edit canonical lessons, then run npm run generate:books.
% canonical-source-hash: fnv1a64:b6d4fb94af652318
% canonical-lessons: FR-C506-match-nul, FR-C506-classement, FR-C506-eliminer, FR-C506-qualifier, FR-C506-vainqueur

\chapter{Le résultat}
\label{ch:fr-le-resultat}
\hlchaptermodality{506}

\begin{chapteropening}
\textbf{By the end of this chapter:} \emph{I can talk about results and rankings.}

Complete the last lesson of chapter 506: I can talk about results and rankings.
\end{chapteropening}

\section[le match nul]{le match nul — draw, tie (in sport)}

\label{lesson:FR-C506-match-nul}



\begin{quote}
Before the new one: say the French for pitch, then the French for to score.
\end{quote}

\subsection*{You'll want to know: le match nul}
\textbf{le match nul} — \textquotedblleft{}draw, tie (in sport)\textquotedblright{}.

\begin{quote}
Les deux équipes ont fait match nul.
\end{quote}

\begin{quote}
\emph{The two teams drew.}
\end{quote}

\emph{Faire match nul} — to draw.

\begin{cousinweb}
\emph{le match nul} joins \emph{match}, borrowed from English, and \emph{nul} \textquotedblleft{}null, none\textquotedblright{}, from Latin \emph{nūllus} \textquotedblleft{}not any\textquotedblright{} (Lewis \& Short).
\end{cousinweb}

\begin{grammarlens}[title={What you've built}]
One more word for telling, explaining and describing at B1.
\end{grammarlens}

\subsection*{Your turn}
\begin{itemize}
\raggedright
  \item \emph{Say it:} \emph{le match nul}
  \item \emph{Say it:} \emph{le match nul}, once more
  \item \emph{Say it:} a sentence of your own with \emph{le match nul}
  \item \emph{Recall:} say the French for pitch, then the French for to score, then say \emph{le match nul} again
\end{itemize}

\begin{tcolorbox}[breakable,colback=teal!4,colframe=teal!35!black,title={Before you move on}]
What does le match nul mean? (\textquotedblleft{}Draw, tie (in sport)\textquotedblright{}.) Say it once more.
\end{tcolorbox}
\section[le classement]{le classement — ranking, league table}

\label{lesson:FR-C506-classement}



\begin{quote}
Before the new one: say the French for to score, then the French for draw.
\end{quote}

\subsection*{You'll want to know: le classement}
\textbf{le classement} — \textquotedblleft{}ranking, league table\textquotedblright{}.

\begin{quote}
Notre équipe est troisième au classement.
\end{quote}

\begin{quote}
\emph{Our team is third in the league.}
\end{quote}

Built on \emph{classer} \textquotedblleft{}to sort, to rank\textquotedblright{}.

\begin{cousinweb}
\emph{le classement} is built on \emph{classer} \textquotedblleft{}to rank\textquotedblright{}, from \emph{classe}, a learned borrowing from Latin \emph{classis} \textquotedblleft{}a class, a division\textquotedblright{} (Lewis \& Short).
\end{cousinweb}

\begin{grammarlens}[title={What you've built}]
One more word for telling, explaining and describing at B1.
\end{grammarlens}

\subsection*{Your turn}
\begin{itemize}
\raggedright
  \item \emph{Say it:} \emph{le classement}
  \item \emph{Say it:} \emph{le classement}, once more
  \item \emph{Say it:} a sentence of your own with \emph{le classement}
  \item \emph{Recall:} say the French for to score, then the French for draw, then say \emph{le classement} again
\end{itemize}

\begin{tcolorbox}[breakable,colback=teal!4,colframe=teal!35!black,title={Before you move on}]
What does le classement mean? (\textquotedblleft{}Ranking, league table\textquotedblright{}.) Say it once more.
\end{tcolorbox}
\section[éliminer]{éliminer — to knock out, to eliminate}

\label{lesson:FR-C506-eliminer}



\begin{quote}
Before the new one: say the French for draw, then the French for ranking.
\end{quote}

\subsection*{You'll want to know: éliminer}
\textbf{éliminer} — \textquotedblleft{}to knock out, to eliminate\textquotedblright{}.

\begin{quote}
Notre équipe a été éliminée en demi-finale.
\end{quote}

\begin{quote}
\emph{Our team was knocked out in the semi-final.}
\end{quote}

Passive in sport: \emph{être éliminé}.

\begin{cousinweb}
\emph{éliminer} is a learned borrowing from Latin \emph{ēlīmināre} \textquotedblleft{}to turn out of doors\textquotedblright{} (EtymDB), built on \emph{līmen} \textquotedblleft{}threshold\textquotedblright{} (Lewis \& Short).
\end{cousinweb}

\begin{grammarlens}[title={What you've built}]
One more word for telling, explaining and describing at B1.
\end{grammarlens}

\subsection*{Your turn}
\begin{itemize}
\raggedright
  \item \emph{Say it:} \emph{éliminer}
  \item \emph{Say it:} \emph{éliminer}, once more
  \item \emph{Say it:} a sentence of your own with \emph{éliminer}
  \item \emph{Recall:} say the French for draw, then the French for ranking, then say \emph{éliminer} again
\end{itemize}

\begin{tcolorbox}[breakable,colback=teal!4,colframe=teal!35!black,title={Before you move on}]
What does éliminer mean? (\textquotedblleft{}To knock out, to eliminate\textquotedblright{}.) Say it once more.
\end{tcolorbox}
\section[qualifier]{qualifier — (se qualifier) to qualify; to describe as}

\label{lesson:FR-C506-qualifier}



\begin{quote}
Before the new one: say the French for ranking, then the French for to knock out.
\end{quote}

\subsection*{You'll want to know: qualifier}
\textbf{qualifier} — \textquotedblleft{}(se qualifier) to qualify; to describe as\textquotedblright{}.

\begin{quote}
L'équipe s'est qualifiée pour la finale.
\end{quote}

\begin{quote}
\emph{The team qualified for the final.}
\end{quote}

Reflexive in sport: \emph{se qualifier pour}. Also to describe: \emph{qualifier quelqu'un de menteur}.

\begin{cousinweb}
\emph{qualifier} is a learned borrowing from Medieval Latin \emph{quālificāre}, built from \emph{quālis} \textquotedblleft{}of what kind\textquotedblright{} and \emph{facere} \textquotedblleft{}to make\textquotedblright{} (Lewis \& Short).
\end{cousinweb}

\begin{grammarlens}[title={What you've built}]
One more word for telling, explaining and describing at B1.
\end{grammarlens}

\subsection*{Your turn}
\begin{itemize}
\raggedright
  \item \emph{Say it:} \emph{qualifier}
  \item \emph{Say it:} \emph{qualifier}, once more
  \item \emph{Say it:} a sentence of your own with \emph{qualifier}
  \item \emph{Recall:} say the French for ranking, then the French for to knock out, then say \emph{qualifier} again
\end{itemize}

\begin{tcolorbox}[breakable,colback=teal!4,colframe=teal!35!black,title={Before you move on}]
What does qualifier mean? (\textquotedblleft{}(se qualifier) to qualify; to describe as\textquotedblright{}.) Say it once more.
\end{tcolorbox}
\section[le vainqueur]{le vainqueur — winner, victor}

\label{lesson:FR-C506-vainqueur}



\begin{quote}
Before the new one: say the French for to knock out, then the French for (se qualifier) to qualify.
\end{quote}

\subsection*{You'll want to know: le vainqueur}
\textbf{le vainqueur} — \textquotedblleft{}winner, victor\textquotedblright{}.

\begin{quote}
Le vainqueur du match jouera la finale.
\end{quote}

\begin{quote}
\emph{The winner of the match will play in the final.}
\end{quote}

Traditionally the same form for a woman: \emph{elle est sortie vainqueur}; \emph{la vainqueure} is also found today, especially in Québec.

\begin{cousinweb}
\emph{le vainqueur} is built on \emph{vaincre} \textquotedblleft{}to defeat\textquotedblright{}, which is inherited from Latin \emph{vincere} \textquotedblleft{}to conquer\textquotedblright{} (Lewis \& Short).
\end{cousinweb}

\begin{grammarlens}[title={What you've built}]
One more word for telling, explaining and describing at B1.
\end{grammarlens}

\subsection*{Your turn}
\begin{itemize}
\raggedright
  \item \emph{Say it:} \emph{le vainqueur}
  \item \emph{Say it:} \emph{le vainqueur}, once more
  \item \emph{Say it:} a sentence of your own with \emph{le vainqueur}
  \item \emph{Recall:} say the French for to knock out, then the French for (se qualifier) to qualify, then say \emph{le vainqueur} again
\end{itemize}

\begin{tcolorbox}[breakable,colback=teal!4,colframe=teal!35!black,title={Before you move on}]
What does le vainqueur mean? (\textquotedblleft{}Winner, victor\textquotedblright{}.) Say it once more.
\end{tcolorbox}
